\documentclass[UTF8,a4paper,zihao=-4,fontset=fandol]{ctexart}
\usepackage[left=2.5cm,right=2.5cm,top=2.5cm,bottom=2.5cm]{geometry}
\usepackage{booktabs,longtable,tabularx,array}
\usepackage{amssymb}
\usepackage{hyperref}
\hypersetup{colorlinks=true,linkcolor=black,urlcolor=blue}
\setlength{\parindent}{2em}
\linespread{1.2}
\newcolumntype{L}[1]{>{\raggedright\arraybackslash}p{#1}}

\begin{document}
\pagestyle{plain}
\begin{center}
{\heiti\zihao{2} AI工具使用详情}
\end{center}

\section{工具与使用范围}

本队在 A 题建模过程中使用 OpenAI Codex（GPT-5）作为辅助工具。使用范围包括：题目结构梳理、数值代码调试建议、结果一致性检查、反事实与灵敏度分析脚本草拟、图表排版、LaTeX 排版和中文表达整理。AI 工具不替代参赛队对控制方程、边界条件、经验公式、达标判据和最终结论的决定责任。

\section{主要使用阶段与过程}

\begin{longtable}{L{2.2cm}L{4.0cm}L{5.6cm}}
\toprule
阶段 & 典型提问或任务 & 采纳、修改与核验方式 \\
\midrule
\endfirsthead
\toprule
阶段 & 典型提问或任务 & 采纳、修改与核验方式 \\
\midrule
\endhead
审题与拆解 & 比较 A、C 两题的计算难度；把 A 题四问拆成输入、模型、输出和检验任务。 & 仅采纳结构化拆解；题目数值、单位、附件和输出格式逐项回到原题 PDF 与 Excel 核对。 \\
模型建立 & 给出圆柱径向热传导、水分扩散、第三类边界及收缩坐标的候选写法。 & 参赛队需检查符号方向、量纲和经验式；问题 4 明确补充“均匀径向收缩、物质坐标”假设，并把移动域守恒不足列为限制。 \\
代码实现 & 辅助编写有限体积、BDF、阈值事件、Excel 导出与输入/输出验证脚本。 & 使用两级或三级网格重算；检查非负性、固定域通量平衡、时间/空间覆盖、Excel 公式错误和末行阈值。所有正式结果由本地脚本重新运行产生。 \\
结果分析 & 提醒问题 3、4 同时改变物性关联式与半径，建议增加 $2\times2$ 反事实和单因素灵敏度。 & 采纳分析框架；数值由独立脚本计算。论文将反事实限定为模型内条件比较，不作实验因果解释。 \\
图表与写作 & 生成中文图件草稿、表格排版、论文结构和文字初稿。 & 图表数据仅来自正式 CSV/JSON；正文数值与机器结果逐项比对；删除不能由题目、代码或结果支持的表述。 \\
\bottomrule
\end{longtable}

\section{关键人工核验清单}

正式提交前，三名队员应共同完成并保留以下核验记录：
\begin{enumerate}
  \item 核对题目四问、附录 2--4 公式、所有单位和初始条件；
  \item 从项目根目录运行正式流程，并确认 $63$ 项自动检查没有失败项；
  \item 抽查表 1--表 6 与 `result1.xlsx`--`result4.xlsx` 的对应单元格；
  \item 复核问题 3 的 BDF 事件根 $57.1686\,\mathrm{h}$、独立全隐式解 $57.2108\,\mathrm{h}$ 与安全提交末行 $57.2500\,\mathrm{h}$ 的定义差别；
  \item 复核问题 4 的 BDF 事件根 $50.8243\,\mathrm{h}$、独立全隐式解 $50.8634\,\mathrm{h}$ 与安全提交末行 $50.8833\,\mathrm{h}$ 的定义差别；
  \item 确认问题 3、4 的直接差异未被表述为纯收缩效应；
  \item 检查长期炉况延拓、端部忽略、潜热闭合不足和移动域闭合均已在论文中披露；
  \item 逐页检查最终 PDF 的匿名性、公式、中文字体、图表、页码和引用；
  \item 由全体队员确认 AI 生成或修改的内容均已理解，并对最终提交负责。
\end{enumerate}

\section{未采纳或降级使用的内容}

未采纳把问题 3 与问题 4 的时长差直接归因于收缩的表述，因为两问还使用了不同物性经验式。未把长期空气水分降低导致问题 3 计算时间大幅增加解释为真实工艺规律，而将其降级为经验扩散式低含水率外推风险。未宣称问题 4 已完成严格的移动域局部质量守恒验证，因为题目没有提供固体速度、孔隙率和干物质密度。

\section{可复现性说明}

核心求解、反事实、灵敏度、图件与 Excel 均由项目中的版本化脚本生成。正式运行使用 Python $3.14.3$、NumPy $2.4.6$、Pandas $3.0.3$、SciPy $1.17.1$ 和 Matplotlib $3.10.9$。原始数据只读，输出分别保存于 `outputs/results/`、`outputs/tables/`、`outputs/figures/` 和 `outputs/submissions/problem\_a/`。

\end{document}
